% !TeX root = ../tesis.tex

El objetivo de este capitulo es aislar y formalizar la capa de certificacion sobre la que se apoya la arquitectura del \zkBridge. La decisión central ya fue anticipada en el Capítulo 5: en lugar de demostrar dentro de un circuito \ZK la validez completa del consenso nativo de \Cardano, usamos \Mithril como mecanismo intermedio de autenticación de estado \cite{cryptoeprint:2021/916,MithrilDocsOverview}.

Como motivación, observamos que por un lado, la validación directa de \Praos traería al \zkBridge detalles de alto costo y fuerte inestabilidad semántica, como \VRF, firmas \KES, nonces de epoca, seleccion de lideres y reglas de cadena. Por otro lado, \Mithril expone una interfaz más adecuada para verificación compacta: registros cerrados de participantes, \textit{aggregate verification keys} (\AVK), cadenas de certificados, y mensajes de protocolo con estructura fija. Usaremos a \Mithril como una abstracción criptográfica del consenso de \Cardano suficientemente fuerte para autenticar el estado relevante y suficientemente compacta para ser reutilizada por la capa de aplicación del \zkBridge.

\Mithril no sustituye al consenso de \Cardano en sentido pleno: no decide qué cadena es válida ni produce bloques. Hace el trabajo de un \LightClient: certificar ciertos datos derivados del estado de \Cardano, de manera que un sistema externo pueda tratarlos como objetos autenticados. Los artefactos relevantes para este proyecto son los certificados \StakeDistribution y \TransactionSnapshot \cite{MithrilDocsCertChain,MithrilDocsClient}.

La capa de actualización mantiene en \ChainB un \UTxO canónico de \StakeDistribution, renovado mediante \StakeDistTx, cuyo \textit{datum} contiene el \StakeDistributionCertificate vigente. La capa de aplicación consume como evidencia un \TransactionSnapshotCertificate y una prueba de inclusion de la \LockTx correspondiente en dicho \TransactionSnapshotCertificate.

El Capítulo 6 ya habia utilizado esta idea, pero de forma deliberadamente operacional. Lo que faltaba justificar era por qué el \UTxO de \StakeDistribution debe considerarse auténtico, cómo se encadena con el estado anterior del \UTxO, qué informacion comprometela \AVK, qué significa que un certificado de \Mithril esté correctamente formado y de qué manera dicha información llega hasta los argumentos \ZK verificados.

No vamos a hablar sobre \textbf{todo} el protocolo \Mithril. Primero describimos el protocolo \STM y sus esquemas de firma, luego formalizamos la cadena de certificados, vemos qué certificados son relevantes al \zkBridge y la interfaz expuesta por el \MithrilAggregator de \Mithril. Sobre esa base, formulamos la relacion $\NP$ del circuito \STM y vemos su implementación en \HaloTwo con \KZG. Una vez fijado todo esto, vamos a saber cómo llegar a implementar el \zkBridge por completo, en particular con los argumentos \ZK relevantes, lo cual desarrollamos en el Capítulo 8.
